\section{Tính hệ thống của khoa học}

\

Khi tổng hợp đủ nhiều các khẳng định và giả thuyết đã được kiểm chứng, chúng ta sẽ có thể rút ra được những tính chất chung, hay còn gọi là các \defText{quy luật}. Những quy luật mà bạn đọc có thể đã biết đến bao gồm:
\begin{itemize}
    \item \textcolor{colorEmphasisCyan}{Định luật vạn vật hấp dẫn}: Định luật này giải thích quy luật tương tác giữa các vật thể có khối lượng trong vũ trụ. Niu-tơn\footnote{Isaac Newton (1643 -- 1727).} không bỗng nhiên nghĩ ra định luật này chỉ vì một quả táo rơi. Nó là kết quả của sự tổng hợp dữ liệu khổng lồ trước đó. Kép-lơ\footnote{Johannes Kepler (1571 -- 1630).} đã tổng hợp dữ liệu về chuyển động của các hành tinh quan sát bởi Bra-hê\footnote{Tycho Brahe (1546 -- 1601).} để rút ra được ba định luật chuyển động hành tinh, từ đó Niu-tơn mới có thể sử dụng toán học để thống nhất hiện tượng ``vật rơi dưới đất'' và ``hành tinh quay trên trời'' thành một công thức duy nhất. Định luật này đã được kiểm chứng qua vô số thực nghiệm và quan sát, từ việc dự đoán quỹ đạo của các hành tinh đến việc giải thích hiện tượng thủy triều trên Trái Đất.

    \item \textcolor{colorEmphasis}{Định luật Men-đen về di truyền học}\footnote{Gregor Mendel (1822 -- 1884).}: Định luật này giải thích cách thức di truyền các đặc điểm từ thế hệ này sang thế hệ khác, dựa trên các thí nghiệm với cây đậu Hà Lan, từ đó tạo nên nền tảng của di truyền học hiện đại. Trước Men-đen, người ta tin rằng đặc điểm của cha mẹ sẽ ``hòa quyện'' vào nhau (như pha màu sơn) để tạo ra con cái. Men-đen nghi ngờ điều này khi thấy những thế hệ con cháu có những đặc điểm không giống với cha mẹ, như có những cây đậu thân thấp thỉnh thoảng vẫn xuất hiện từ cặp bố mẹ thân cao. Ông đã lai tạo hơn $29000$ cây đậu Hà Lan trong suốt $8$ năm. Trong quá trình nghiên cứu, ông tập trung vào các tính trạng đơn lẻ (màu hạt, độ cao thân). Sau khi tổng hợp dữ liệu tính trạng của nhiều thế hệ (cũng có thể hiểu là rất nhiều khẳng định đơn lẻ) và kết hợp chúng với những phát hiện trong thống kê, ông đi đến kết luận rằng di truyền không phải là sự hòa tan, mà là sự truyền lại của các \emph{đơn vị riêng biệt} (mà sau này ta gọi là \emph{gien}).

    \item \textcolor{colorEmphasisGreen}{Định luật bảo toàn và chuyển hóa năng lượng}: Một trong những quy luật cơ bản nhất, khẳng định năng lượng không tự nhiên sinh ra cũng không mất đi. Tuy quy luật này thường được phát biểu trong phạm trù của vật lí, nó không xuất phát từ một thí nghiệm duy nhất mà là sự giao thoa của nhiều ngành. Giun\footnote{James Prescott Joule (1818 -- 1889).} đã dùng một quả nặng rơi để làm quay các cánh quạt trong nước và chứng minh rằng công cơ học biến thành nhiệt năng theo một tỉ lệ không đổi. Mai-ơ\footnote{Julius Robert von Mayer (1814 -- 1878).}, sau khi quan sát được sự khác biệt của màu máu giữa các vùng, ông đã đưa ra kết luận rằng thức ăn không chỉ tạo ra nhiệt lượng mà còn tạo ra lực chuyển động và nhiệt giữ ấm cho cơ thể, đóng góp vào việc xây dựng định luật này.

          Định luật bảo toàn và chuyển hóa năng lượng cũng là một trong những nguyên nhân khách quan dẫn đến sự tồn tại của chủ nghĩa duy vật biện chứng của Các Mác\footnote{Karl Marx (1818 -- 1883).} và Ăng-ghen\footnote{Friedrich Engels (1820 -- 1895).}, khi họ sử dụng quy luật này để giải thích sự vận động luân hồi của vật chất.
\end{itemize}

Bạn đọc có thể thấy rằng các định luật tưởng chừng như đơn giản đó đã phải dựa trên nền tảng dữ liệu khổng lồ và vô số thực nghiệm và quan sát để có thể được phát hiện ra. Những gì chúng ta biết hôm nay là kết quả của hàng ngàn năm quan sát và thử sai của các thế hệ đi trước. Khoa học không đứng yên mà nó luôn tự cập nhật, sửa đổi và hoàn thiện khi con người phát hiện ra những dữ liệu mới chính xác hơn.
